\documentclass[11pt]{letter}
\usepackage[a4paper,margin=1in]{geometry}
\usepackage{hyperref}

\signature{Sanjiang Li \\
  Centre for Quantum Software and Information \\
  University of Technology Sydney, Australia \\
  %ORCID: 0000-0002-3332-2546 \\
  E-mail: sanjiang.li@uts.edu.au}

%\address{Sanjiang Li \\
%  Centre for Quantum Software and Information \\
%  University of Technology Sydney \\
%  PO Box 123, Broadway NSW 2007, Australia}

\begin{document}

\begin{letter}{}
%The Editor-in-Chief \\
%  IEEE Transactions on Computer-Aided Design of Integrated Circuits and Systems \\
 % IEEE Council on Electronic Design Automation}

\opening{Dear Editor,}

I am pleased to submit the enclosed manuscript, ``\emph{dSABRE: A
SABRE-Style Router for Multi-Core Distributed Quantum Computers},''
for consideration as a Regular Paper in the IEEE Transactions on
Computer-Aided Design of Integrated Circuits and Systems (TCAD).

%\medskip
\noindent\textbf{Scope and contribution.}
The paper addresses the routing problem for circuits compiled to
multi-core distributed quantum computers (DQCs), where minimising
EPR-pair consumption is the dominant compilation objective.  We
present dSABRE, a SABRE-style router whose inter-core teleportation
scorer combines a gate-centric five-term cost (staging, capacity,
front-layer gain, hop direction, and decay-weighted lookahead) with
two architectural mechanisms --- a proactive congestion-relief pass
and a BFS-layer inter-core extended set.  On 18 MQT-Bench circuits
at 25, 36, and 64 logical qubits, dSABRE reduces geometric-mean EPR
consumption by 41--44\% over TeleSABRE (the prior state-of-the-art
SABRE-style distributed router) and by 16--68\% over pytket-dqc, and
scales to 360-qubit QFT circuits on a 486-physical-qubit
architecture.  The contribution fits squarely within TCAD's scope on
algorithms and methods for automated design of integrated quantum
systems.

%\medskip
\noindent\textbf{Originality and prior versions.}
This submission is original work.  No earlier peer-reviewed version
of this paper has been published or accepted for publication at any
journal, conference, or workshop, and no portion of the text overlaps
with prior peer-reviewed publications by the author.

%\medskip
\noindent\textbf{Concurrent submission.}
The manuscript is not under review at any other journal, conference,
or workshop, and will not be submitted elsewhere during the TCAD
review process.

%\medskip
\noindent\textbf{Open access.}
I am submitting under the \emph{Traditional} (non-Open-Access)
publication option.

%\medskip
\noindent\textbf{Generative AI disclosure.}
In accordance with TCAD's policy on generative AI tools (§G), the
acknowledgments section discloses the use of Anthropic's Claude
(via Claude Code) for: refactoring portions of the dSABRE
implementation, authoring benchmark scripts and aggregating
per-circuit JSON results into tables, drafting LaTeX sources for two
figures, and copy-editing prose.  All AI-generated content was
reviewed, verified, and corrected by the author, who takes full
responsibility for correctness and originality.

%\medskip
\noindent\textbf{Conflicts of interest.}
The author declares no conflicts of interest with any current TCAD
Associate Editor.

%\medskip
\noindent\textbf{Artefact availability.}
The implementation, benchmark suites, per-circuit results, and
online appendices referenced in the paper are publicly available at
\url{https://github.com/ebony72/dsabre}.

%\medskip
Thank you for considering this submission.  I look forward to the
reviewers' feedback.

\closing{Sincerely,}

\end{letter}
\end{document}
